\section{Host rigidity from the literal adjoint and local deformation conditions}
\label{sec:host-rigidity-v42}

The remaining question called \texttt{HOST\_RIGID\_S} is a
reconstruction matching question: do the particular deformations in
the source give Bloch--Kato classes for the particular polarized
adjoint to which ADK6 applies?  This section answers that question.
The pairing is the source's antidiagonal matrix, the scalar twist is
exactly $\epsilon^{-1}$, and the local argument includes arbitrary
dual-number changes of basis and unramified diagonal variations.
The conclusion does not construct the automorphic points or prove
the separate quantitative host-number bound.

\subsection{The original hypotheses, in ordinary notation}

Fix an algebraic closure $\Omega$ of $\mathbb Q$, an odd prime
$\ell$, and a finite $\mathbb Z_\ell$-subalgebra
$O\subset\overline{\mathbb Q}_\ell$ which is a discrete valuation
ring.  Let $N\ne0$ be even, let
$\psi:(\mathbb Z/N\mathbb Z)^\times\to O^\times$ be a Dirichlet
character, and let $e,o\ge1$.  For a prime $p$, write $\psi_p$ for
the CRT component of $\psi$ supported at the $p$-part of the modulus.
This is the source's \texttt{comp p}; it is different from the
decomposition according to the order of the character.  Write
$\psi=\theta\kappa$, where $\theta$ has order prime to $\ell$
and $\kappa$ has $\ell$-power order.  These are
\texttt{th0} and \texttt{kap}.  Put
\[
 c_2=v_2(\operatorname{cond}\psi_2),\qquad
 S_b=\{p\mid N:p\ne\ell,\ \kappa_p\ne1\},\qquad
 T_b=\{\ell\}\cup S_b.
 \tag{HSR.1}
\]
The complete \texttt{HostData} assumptions are
\[
 \psi(-1)=(-1)^e,\quad c_2\ge4,\quad
 \theta_p^{\,2o}=1\ (p\ne2),\quad
 8op^2<2^{c_2}\ (p\in T_b).
 \tag{HSR.2}
\]
The prime-to-$\ell$ and $\ell$-power order parts exist because
$\psi$ is a character of a finite group.  The explicit exponent
in the source's definition of \texttt{lPart} is the CRT projector
on the $\ell$-primary subgroup of the finite cyclic group generated
by $\psi$.

Let $S$ be exactly the finite primes dividing $\ell N$, let
$\Gamma=G_{\mathbb Q,S}$, and let
$\pi:G_{\mathbb Q}\twoheadrightarrow\Gamma$ be the specified
quotient.  For each $p\in S$ use the stored valuation above $p$,
with decomposition and inertia groups $D_p,I_p\subset G_{\mathbb Q}$.
The character $\eta:\Gamma\to\{\pm1\}$ and the character
$\Psi:\Gamma\to O^\times$ satisfy
\[
 \eta\pi=\eta_4,\qquad
 \Psi\pi=\epsilon^e\psi_{\mathrm{Gal}}.
 \tag{HSR.3}
\]
Here $\eta_4$ is the quadratic character of $\mathbb Q(i)$,
and $\epsilon$ is the actual $\ell$-adic cyclotomic character.
A stored element $c\in\Gamma$ is the image of an actual complex
conjugation $c_\infty\in G_{\mathbb Q}$.  This is precisely the
content of \texttt{IsHostFrame}, including its choices of valuations,
quotient map, distinguished $\ell$-place, and subset $S_b$.
An actual complex conjugation sends $i$ to $-i$, so
$\eta(c)=\eta_4(c_\infty)=-1$.

Let
\[
 J=\begin{pmatrix}0&0&1\\0&1&0\\1&0&0\end{pmatrix},\qquad
 \vartheta(a)=Ja^{-t}J,
 \qquad \mathcal G_3=\operatorname{GL}_3\rtimes\{\pm1\},
 \tag{HSR.4}
\]
where $-1$ acts by $\vartheta$.  The source's group
\texttt{GG 3} is exactly $\mathcal G_3$.  If
$R(g)=(a(g),\eta(g))$, its adjoint action on the full space
$V=M_3(\overline{\mathbb Q}_\ell)$ is
\[
 \operatorname{Ad}_R(g)X=
 \begin{cases}
 a(g)Xa(g)^{-1},&\eta(g)=1,\\
 -a(g)JX^tJa(g)^{-1},&\eta(g)=-1.
 \end{cases}
 \tag{HSR.5}
\]

Fix the marked element $g_1$ required by the source and a point
$r:\Gamma\to\mathcal G_3(\overline{\mathbb Q}_\ell)$ satisfying
\texttt{IsCharZeroPt}.  Thus there are a finite $O$-subalgebra
$A\subset\overline{\mathbb Q}_\ell$ which is a DVR, a unit
$t\in\overline{\mathbb Q}_\ell^\times$, and an integral point
$r_A:\Gamma\to\mathcal G_3(A)$ such that
\[
 r_A=\operatorname{conj}_{T_t}r,
 \qquad T_t=\operatorname{diag}(t,1,t^{-1}).
 \tag{HSR.6}
\]
Scalar inclusions in (HSR.6) are understood.  The integral point has
multiplier $\eta$, has continuous matrix entries modulo every power
of $\mathfrak m_A$, has $r_A(c)^{\rm mat}=1$, and satisfies
\[
 r_A(g)^{\rm mat}\equiv
 \operatorname{diag}(\Psi(g),1,\Psi(g)^{-1})\pmod{\mathfrak m_A}
 \qquad(g\in\Gamma).
 \tag{HSR.7}
\]
The remaining normalization and local conditions on this integral
point are retained from the source, although the proof below does
not need to use each of them separately.

For clarity, here is the full conclusion \texttt{Rigid r} to be
proved.  Let $z:\Gamma\to V$ be any continuous cocycle for (HSR.5),
and let $r_\delta$ be any deformation over
$\overline{\mathbb Q}_\ell[\delta]/(\delta^2)$ with the same
multiplier and
\[
 r_\delta(g)^{\rm mat}=(1+\delta z(g))r(g)^{\rm mat}.
 \tag{HSR.8}
\]
For $p\notin T_b$ assume that $z$ vanishes on
\[
 N_p=I_p\cap\ker\eta_4\cap\ker\overline{\Psi\pi}.
 \tag{HSR.9}
\]
At $\ell$, assume that some dual-number basis $h$ makes $r_\delta$
upper triangular on $D_\ell\cap\ker\eta_4$, with inertia diagonal
\[
                  (\Psi,1,\Psi^{-1}).
 \tag{HSR.10}
\]
The right side of (HSR.10) has no $\delta$ term.  At every
$p\in S_b$, assume that such a basis makes the restriction diagonal
on $D_p\cap\ker\eta_4$, again with (HSR.10) on inertia.  These
are literally \texttt{TriIn} and \texttt{DiagIn}; their bases are
not required to be constant or close to the identity.  We must find
one $X\in V$ with
\[
                  z(g)=X-\operatorname{Ad}_r(g)X
                  \qquad(g\in\Gamma).
 \tag{HSR.11}
\]
The sign in (HSR.11) is the source's \texttt{cob} convention.

\subsection{Integral points and the local separating element}

% BEGIN HSR1
\begin{lemma}[HSR1: continuity, normalization, and finite coefficients]
\label{claim:HSR1}
The point $r$ has continuous matrix entries, has multiplier $\eta$,
and satisfies $r(c)=(1,-1)$.  Its adjoint action is continuous.
For every cocycle and deformation in (HSR.8), all matrix coefficients
of $r,z,r_\delta$, together with any fixed finite set of dual-number
bases, lie in one finite extension of $\mathbb Q_\ell$.
\end{lemma}
\begin{proof}
Every element of $A$ is integral over $\mathbb Z_\ell$, hence has
absolute value at most one: otherwise the leading term of a monic
integral equation would have strictly larger absolute value than
all its other terms.  Every unit of $A$ has absolute value one,
by applying the same assertion to it and its inverse.  If $\varpi$
is a uniformizer, write $\ell=u\varpi^b$.  Here $b\ge1$, since
$\ell$ has absolute value less than one and cannot be a unit of
$A$.  Thus $q=|\varpi|=\ell^{-1/b}<1$, and an element of
$\mathfrak m_A^n$ has absolute value at most $q^n$.  Local constancy
modulo $\mathfrak m_A^n$ for every $n$ now proves norm continuity
of every entry of $r_A$, hence of $r$ under (HSR.6).

Conjugation does not change the multiplier.  Also
$\vartheta(T_t)=T_t$.  Since $\eta(c)=-1$, the matrix part of
$T_t r(c)T_t^{-1}$ is consequently
$T_t r(c)^{\rm mat}T_t^{-1}$.  Its equality to $1$ proves
$r(c)^{\rm mat}=1$.  The character $\eta$ is locally constant:
its composition with the quotient map is the continuous finite
character $\eta_4$, and the quotient topology gives continuity of
$\eta$.  Formula (HSR.5), on its two open cosets, proves continuity
of the adjoint action.

The fraction field of $A$ is finite over $\mathbb Q_\ell$ because
$A$ is finite over $O$ and $O$ is finite over $\mathbb Z_\ell$.
Adjoining $t$ gives a finite field containing every entry of $r$ and
its adjoint action.  ADK1's cocycle argument then puts all values of
the chosen $z$ in a finite extension.  Formula (HSR.8) includes the
deformation, and adjoining the finitely many entries of the chosen
bases includes those bases and their inverses.

For the more general \texttt{RigDesc} interface one can start
directly with a continuous dual-number representation on any compact
group.  Take an increasing countable cofinal family $E_i$ of finite
local coefficient fields, as in NR19 and ADK1.  The sets where all
constant and $\delta$ matrix coefficients belong to $E_i$ are
closed subgroups and cover the group.  They are subgroups because
the semidirect-product operations are defined over $E_i[\delta]$;
a matrix over this ring which is invertible over the larger ring
is already invertible over it, since its constant determinant is
nonzero.  Baire's theorem makes one subgroup open.  Adjoining the
coefficients at its finitely many coset representatives and of the
specified basis gives the asserted finite field.  The topology on
each finite field is its subspace topology, so all descended matrix
maps remain continuous.  This proves exactly the coefficient
descent used in \texttt{RigDesc}, rather than assuming that arbitrary
compact subsets of the algebraic closure have finite coefficients.
\end{proof}
% END HSR1

% BEGIN HSR2
\begin{lemma}[HSR2: the dyadic conductor forces separation after reduction]
\label{claim:HSR2}
Suppose $u\in(\mathbb Z/N\mathbb Z)^\times$ has $2$-primary
coordinate equal to $\ell$.  Under (HSR.2),
\[
                         \overline{\psi(u)}^{\,4}\ne1.
 \tag{HSR.12}
\]
\end{lemma}
\begin{proof}
Put $s=v_2(\ell^2-1)\ge3$, write $o=2^jm$ with $m$ odd, and
put $a=v_2(N)$ and $d=s+j+1$.  Squaring $1+2^sb$ with $b$ odd
raises the exact $2$-adic valuation of its difference from $1$ by
one when $s\ge2$; raising to an odd power does not change it.
Induction gives
\[
 v_2(\ell^{4o}-1)=d,\qquad
 2^d\le2o(\ell^2-1)<8o\ell^2<2^{c_2}.
 \tag{HSR.13}
\]
Thus $d<c_2\le a$.

For completeness, an element $v=1+2^d b$ with $b$ odd and $d\ge2$
generates the group $1+2^d\mathbb Z/2^a\mathbb Z$.
The same squaring calculation shows
$v_2(v^{2^k}-1)=d+k$, so its order modulo $2^a$ is exactly
$2^{a-d}$, the size of that group.  If
$\psi_2(u)^{4o}=1$, then the character on
$(\mathbb Z/2^a\mathbb Z)^\times$ kills the generator
$\ell^{4o}$ of this group.  It therefore factors through
$(\mathbb Z/2^d\mathbb Z)^\times$.  Its conductor divides $2^d$,
contradicting the definition of $c_2$.  We have proved
$\psi_2(u)^{4o}\ne1$.

Reduction to the residue field kills all $\ell$-power order roots
of unity.  It is injective on prime-to-$\ell$ roots: if $\zeta^b=1$
with $\ell\nmid b$ and $\bar\zeta=1$, then
$1+\zeta+\cdots+\zeta^{b-1}$ is a unit, and
$(\zeta-1)(1+\cdots+\zeta^{b-1})=0$ implies $\zeta=1$.
Since the $2$-primary unit group has $2$-power order and $\ell$ is
odd, $\kappa_2=1$ and $\theta_2=\psi_2$.  The CRT product
decomposition and (HSR.2) give
\[
                   \theta(u)^{4o}=\psi_2(u)^{4o}\ne1.
 \tag{HSR.14}
\]
If $\overline{\psi(u)}^4=1$, reduction of $\kappa(u)$ gives
$\overline{\theta(u)^4}=1$, and injectivity on prime-to-$\ell$
roots gives $\theta(u)^4=1$.  This contradicts (HSR.14).
Notice also that $\kappa_2=1$ proves $2\notin S_b$.
\end{proof}
% END HSR2

% BEGIN HSR3
\begin{lemma}[HSR3: the exact local subgroup, with cyclotomic value one]
\label{claim:HSR3}
For the stored decomposition group at $\ell$, there is an open
normal subgroup
\[
 H=D_\ell\cap\ker\eta_4\cap\ker(\psi_\ell)_{\mathrm{Gal}}
 \tag{HSR.15}
\]
and an element $g\in H$ with
\[
 \Psi\pi(s)=\epsilon(s)^e\quad(s\in H\cap I_\ell),\qquad
 \epsilon(g)=1,\qquad (\Psi\pi(g))^2-1\in O^\times.
 \tag{HSR.16}
\]
In particular this proves the full \texttt{rig\_sub} output, with
its final inequality $|\epsilon(g)-1|<1$ strengthened to equality
$\epsilon(g)=1$.
\end{lemma}
\begin{proof}
Choose an embedding $j:\Omega\hookrightarrow
\overline{\mathbb Q}_\ell$ inducing the stored valuation.  One
construction embeds the completions of finite subextensions at
that valuation into the local algebraic closure and takes a
compatible system of these finite embeddings.  Each finite
compatibility problem has a solution by extension of embeddings;
compactness of the product of the finite nonempty sets gives the
system.  Restriction identifies $G_{\mathbb Q_\ell}$
topologically with $D_\ell$ and identifies their inertia groups.
Here are the details relevant to the chosen model.  The image
$j\Omega$ is the field of elements algebraic over $\mathbb Q$
in the local closure and is dense there, by the rational-polynomial
approximation and conjugate-separation argument used in NR19.
An element of $D_\ell$ preserves the normalized absolute value,
so extends isometrically to the completion.  It fixes the closure
$\mathbb Q_\ell$ of $\mathbb Q$, and consequently preserves
the elements algebraic over $\mathbb Q_\ell$.  This gives the
inverse to restriction.  Restriction is continuous: fixing a finite
set $T\subset\Omega$ pulls back to the open subgroup fixing
$\mathbb Q_\ell(j(T))$.  The continuous bijection from the compact
local Galois group to the Hausdorff group $D_\ell$ is a
homeomorphism.  Its valuation-preserving construction identifies
inertia and preserves the action on every root of unity.

The image of local inertia under $\epsilon$ is all of
$\mathbb Z_\ell^\times$.  Indeed
$\Phi_{\ell^m}(1+T)$ is Eisenstein at $\ell$: its constant term
is $\ell$, its reduction modulo $\ell$ is
$T^{\ell^{m-1}(\ell-1)}$, and it is monic.  Thus
$\mathbb Q_\ell(\mu_{\ell^m})/\mathbb Q_\ell$ is totally
ramified of degree $\ell^{m-1}(\ell-1)$, and inertia maps onto
$(\mathbb Z/\ell^m\mathbb Z)^\times$.  Compactness gives the
surjection on the inverse limit as well.

Take any lift $\sigma_0$ of arithmetic Frobenius in the local
unramified quotient.  Choose $\tau\in I_\ell$ with
$\epsilon(\tau)=\epsilon(\sigma_0)^{-1}$ and set
$\sigma=\sigma_0\tau$.  Then $\epsilon(\sigma)=1$.
Roots of unity of order prime to $\ell$ lie in the maximal
unramified extension: they are the unique Hensel lifts of roots
of the same order in finite residue fields.  Hence $\sigma$
still acts on them by raising to the $\ell$th power.  Its action
on $\mu_N$ therefore has $2$-primary coordinate $\ell$.
HSR2 gives $\overline{\psi_{\mathrm{Gal}}(\sigma)}^4\ne1$.

Since $\epsilon(\sigma)=1$, the element $\sigma$ fixes all
$\ell$-power roots of unity and kills $(\psi_\ell)_{\mathrm{Gal}}$.
Set $g=\sigma^2$.  This kills $\eta_4$, so $g\in H$, and
\[
 \overline{\Psi\pi(g)}^2
   =\overline{\psi_{\mathrm{Gal}}(\sigma)}^4\ne1.
\]
An element of the local ring $O$ is a unit exactly when its residue
is nonzero, proving the last assertion in (HSR.16).
Both characters cutting out $H$ have finite image and open kernel,
so $H$ is open and normal in $D_\ell$.  Finally inertia at $\ell$
fixes all prime-to-$\ell$ roots of unity, so on it
$\psi_{\mathrm{Gal}}=(\psi_\ell)_{\mathrm{Gal}}$.
Equations (HSR.3) and (HSR.15) give the first assertion of
(HSR.16).
\end{proof}
% END HSR3

% BEGIN HSR4
\begin{lemma}[HSR4: separation in every triangular basis]
\label{claim:HSR4}
Suppose $g\in\Gamma$ has $\eta(g)=1$ and
$\Psi(g)^2-1\in O^\times$.  In any basis in which $r(g)^{\rm mat}$
is upper triangular, its diagonal entries $\lambda_1,\lambda_2,
\lambda_3$ satisfy
\[
 |\lambda_i|=1,\qquad |\lambda_i-\lambda_j|=1\ (i\ne j),
 \qquad \lambda_i\ne u\lambda_j\quad(|u-1|<1, i\ne j).
 \tag{HSR.17}
\]
This is the exact \texttt{rig\_eig} conclusion.
\end{lemma}
\begin{proof}
The three entries $d=(\Psi(g),1,\Psi(g)^{-1})$ are units with
pairwise distinct residues.  Indeed $\Psi(g)^2-1$ being a unit
makes both $\Psi(g)-1$ and $\Psi(g)+1$ units, and the other
differences follow after multiplication by $\Psi(g)^{-1}$.
By HSR1 and (HSR.7), the integral model $M=r_A(g)^{\rm mat}$
has entries of absolute value at most one and reduces to
$\operatorname{diag}(\bar d_1,\bar d_2,\bar d_3)$.
Its characteristic polynomial is monic integral.  Every root has
absolute value at most one, by the same leading-term argument as
in HSR1.  Reducing its factorization into linear factors gives
\[
 \prod_i(T-\bar\lambda_i)=\prod_i(T-\bar d_i).
\]
The roots on the right are distinct and nonzero.  Unique
factorization over the residue algebraic closure makes the
$\bar\lambda_i$ a permutation of them.  This proves the first
two assertions of (HSR.17).  The triangular diagonals are these
roots because $\eta(g)=1$ makes (HSR.6) ordinary matrix
conjugation.  Finally $|\lambda_j(u-1)|<1$ and
$|\lambda_i-\lambda_j|=1$ imply
$|\lambda_i-u\lambda_j|=1$ by the ultrametric inequality.
\end{proof}
% END HSR4

\subsection{The local deformation is a finite adjoint class}

% BEGIN HSR5
\begin{proposition}[HSR5: the precise RigFin and RigLoc implication]
\label{claim:HSR5}
Let $D$ be an open subgroup of a local absolute Galois group at
$\ell$, let $r:D\to\mathcal G_3(\overline{\mathbb Q}_\ell)$
have locally constant multiplier and continuous matrix entries,
let $z$ be a continuous adjoint cocycle, and let $r_\delta$ be a
deformation with the same multiplier satisfying (HSR.8).
Equivalently, $D$ may be an abstract topological group equipped
with the specified open embedding into $G_{\mathbb Q_\ell}$;
transport along that embedding gives the present formulation.
Suppose an open normal
$H\subset D\cap\ker\eta$ and a dual-number basis give an upper
triangular representation on $H$ whose diagonal characters
$\kappa_1,\kappa_2,\kappa_3$ satisfy
\[
 \kappa_i|_{H\cap I}=\epsilon^{w_i},\qquad
 w_1>w_2>w_3,
 \qquad \bar\kappa_i\ne\bar\kappa_j\epsilon\quad(i<j).
 \tag{HSR.18}
\]
The bar here means reduction modulo $\delta$, not modulo $\ell$.
Then the original adjoint cocycle has a primitive in
$M_3(\overline{\mathbb Q}_\ell)\otimes_{\mathbb Q_\ell}
B_{\mathrm{cris}}$ on all of $D$.  This conclusion is exactly
the source's \texttt{finiteCocycles}, including its tensor order
and its diagonal Galois action.
\end{proposition}
\begin{proof}
The coefficient descent in HSR1 gives a finite field $E$ containing
the representation, cocycle, and basis.  It applies here because
open local Galois subgroups are compact.  All polynomial equations,
the upper-triangular conditions, and the non-equalities in
(HSR.18) descend along the injective inclusion
$E[\delta]\hookrightarrow\overline{\mathbb Q}_\ell[\delta]$.
On $H$, both $r$ and $r_\delta$ are ordinary linear
representations, because their multiplier is one.  Under the
identification $H=G_L$ for a finite extension of the local base,
$H\cap I$ is precisely $I_L$.

We recall the actual calculation from RX2 to record why all these
hypotheses, including the dual-number ones, suffice.  For
$A=E$ or $E[\delta]$, an upper triangular representation with
(HSR.18) is crystalline as an underlying $E$-representation.
Each $\kappa_i\epsilon^{-w_i}$ is unramified.  Even when it
contains a nontrivial unramified $\delta$-variation, its underlying
$E$-representation is unramified and hence crystalline; twisting
by $\epsilon^{w_i}$ preserves this property.  For the induction
step, let $V_0$ be the first two (or first $n-1$) coordinates and
put $W=V_0\otimes_A A(\kappa_n^{-1})$.  By induction and tensor
and quotient stability, $W$ is crystalline.  Its $E$-character
constituents are $\bar\kappa_i/\bar\kappa_n$, with repetitions
when $A$ has dual numbers.  All inertia exponents are positive,
so $\operatorname{Fil}^0D_{\mathrm{dR}}(W)=0$ in the convention
where $\epsilon$ has Hodge--Tate weight $-1$.
The constituents of $W^\vee(1)$ are nontrivial by (HSR.18),
whence $H^0(L,W^\vee(1))=0$.

The local Euler characteristic and Tate duality give
\[
 H^2(L,W)=0,\qquad
 \dim_E H^1(L,W)=\dim_E H^0(L,W)+[L:\mathbb Q_\ell]\dim_EW.
\]
The Bloch--Kato finite dimension formula gives the same expression
for $\dim_EH^1_f(L,W)$, since $\operatorname{Fil}^0=0$.
Thus $H^1=H^1_f$ for this $W$.  RX2 supplies the lattice limit
argument deriving these rational statements from the already
installed finite local duality and Euler identities, and BX1--BX4
supply the finite dimension formula for the actual period model.
The last column of the upper triangular representation, divided
by its last diagonal character, is a cocycle in $W$.  Its
$B_{\mathrm{cris}}$ primitive splits the corresponding underlying
$E$-extension after tensoring.  The crystalline comparison
isomorphisms for its submodule and quotient therefore give one
for the extension.  This proves the induction, with all
unramified dual-number diagonal variations included.

Consequently the underlying $E$-self-extension of $r|_H$ defined
by $r_\delta|_H$ is crystalline.  Exactness of $D_{\mathrm{cris}}$
and the comparison isomorphism give an equivariant splitting after
tensoring with $B_{\mathrm{cris}}$.  Choose it compatibly with $E$
by splitting first over the semisimple algebra
$L_0\otimes_{\mathbb Q_\ell}E$.  In the vector-space splitting
already used in (HSR.8), this is precisely an element $b$ with
\[
 z(h)\otimes1=\operatorname{Ad}_{r,\mathrm{cris}}(h)b-b
                  \qquad(h\in H).
 \tag{HSR.19}
\]

Here is the descent to $D$, also valid if $\ell$ divides $[D:H]$.
Subtract the coboundary in (HSR.19) from $z\otimes1$, obtaining
a cocycle $f$ which vanishes on $H$.  Normality and the cocycle
identity show that it factors through $Q=D/H$ with values in
the $H$-invariants.  For $S_f=\sum_{q\in Q}f(q)$, permutation
of the finite sum gives
\[
                  gS_f-S_f=-|Q|f(g).
\]
Thus $b-|Q|^{-1}S_f$ is a primitive on $D$.  The action here is
the full disconnected adjoint on the first tensor factor and
the actual Galois action on $B_{\mathrm{cris}}$ on the second;
no connectedness assumption on $D$ was made.  Extending
coefficients to the algebraic closure carries this primitive to
the required one.  BX1--BX2 identify the period ring and this
literal tensor-primitive condition with \texttt{finiteCocycles}.
Equivariance of scalar extension is immediate on pure tensors
$X\otimes b$ and hence on their sums.  This proves both the
finite-coefficient \texttt{RigFin} and the algebraic-closure
\texttt{RigLoc} interfaces.
\end{proof}
% END HSR5

% BEGIN HSR6
\begin{lemma}[HSR6: the exact inertia primitives away from ell]
\label{claim:HSR6}
Every deformation satisfying (HSR.8)--(HSR.10) has an adjoint
coboundary on inertia at each finite $p\ne\ell$.
\end{lemma}
\begin{proof}
If $p\notin S$, the quotient $\pi$ kills inertia, and
$z(1)=0$, so the assertion holds with primitive zero.
Suppose $p\in S_b$.  By HSR2, $p\ne2$; hence $\eta_4$ is
trivial on $I_p$, since the fourth roots of unity are unramified
at odd primes.  Write the supplied diagonalizing basis as
$h=h_0+\delta h_1$, with $h_0$ invertible, and write
$C=h^{-1}r_\delta(s)h$ for $s\in I_p$.  Its matrix is diagonal
with the fixed inertia entries (HSR.10), so has no $\delta$
coefficient.  Put $R=r(s)^{\rm mat}$.  The constant and
$\delta$ terms of $hC=(1+\delta z(s))Rh$ give
\[
 h_0C_0=Rh_0,\qquad
 h_1C_0=Rh_1+z(s)Rh_0.
\]
With $U=h_1h_0^{-1}$, this says
\[
 z(s)=U-RUR^{-1}
     =\operatorname{Ad}_r(s)(-U)-(-U).
 \tag{HSR.20}
\]
The single matrix $-U$ works for all inertia elements.  This is
the source's \texttt{rig\_diag}, with its sign explicitly fixed.

If $p\in S\setminus T_b$, the cocycle vanishes on $N_p$ from
(HSR.9).  This is a normal subgroup of finite index in $I_p$:
it is the intersection of the kernels of two finite characters,
$\eta_4$ and $\bar\Psi$, with inertia.  The same finite quotient
sum as in HSR5, now without a period-ring tensor, proves that
$z|_{I_p}$ is a coboundary.  In detail $z$ factors through
$I_p/N_p$ with values invariant under $N_p$, and
$-\sum_{q\in I_p/N_p}z(q)/[I_p:N_p]$ is a primitive.  This
also handles $p=2$, without asserting that $\eta_4$ is
unramified there.  These cases exhaust all $p\ne\ell$.
\end{proof}
% END HSR6

% BEGIN HSR7
\begin{proposition}[HSR7: all source tangent classes satisfy literal IsBK]
\label{claim:HSR7}
Let $R=r\pi$ and $Z=z\pi$.  For every allowed tangent in the
source's \texttt{Rigid} definition,
\[
       Z\in Z^1(G_{\mathbb Q},\operatorname{Ad}_R),
       \qquad \texttt{IsBK}_{\ell}(\operatorname{Ad}_R,Z).
 \tag{HSR.21}
\]
\end{proposition}
\begin{proof}
Continuity follows from HSR1 and continuity of the quotient map;
the cocycle law follows by composition.  The away condition is
HSR6, using the stored valuation when $p\in S$ and any valuation
when $p\notin S$.

At $\ell$, use the subgroup $H$ and element $g$ of HSR3 and the
dual-number upper-triangular basis supplied in (HSR.10).  On $H$
the diagonal characters have inertia exponents
\[
                         (w_1,w_2,w_3)=(e,0,-e),
 \tag{HSR.22}
\]
which are strictly decreasing because $e\ge1$.  Reduction of
the basis modulo $\delta$ is an ordinary invertible matrix and
triangularizes $r(g)^{\rm mat}$.  HSR4 gives pairwise distinct
diagonal entries at $g$.  Since $\epsilon(g)=1$, this proves
$\bar\kappa_i(g)\ne\bar\kappa_j(g)\epsilon(g)$ for every
$i<j$, hence the character non-equalities in (HSR.18).
HSR5 gives the required tensor primitive on the entire stored
$D_\ell$.

To check the exact \texttt{IsBK.above} interface, take the embedding
$j$ of HSR3.  A field homomorphism
$\tau:\mathbb Q\to\overline{\mathbb Q}_\ell$ is the unique
rational inclusion, so $j|_{\mathbb Q}=\tau$ as required.
The source's \texttt{decJ} consists of pairs $(s,u)$ such that
$u j=j s$.  It is the graph of restriction from
$G_{\mathbb Q_\ell}$; its first projection has image exactly
the stored $D_\ell$ and its second projection is a topological
isomorphism.  HSR3 supplies the identification with that stored
group, and ADK5 gives the graph topology directly by open
stabilizers of finite global sets.  Transport (HSR.19) along these
maps.  On a pure tensor the transported action is
\[
            X\otimes b\longmapsto
            \operatorname{Ad}_R(s)X\otimes u(b),
\]
which is exactly the source's \texttt{actC}.  Thus the transported
primitive proves its \texttt{finiteCocycles} assertion, not just
a condition on an unspecified local representation.  This
verifies both clauses of the literal \texttt{IsBK}.
\end{proof}
% END HSR7

\subsection{The exact global polarization and the endpoint}

% BEGIN HSR8
\begin{proposition}[HSR8: the automorphic twist and full disconnected adjoint]
\label{claim:HSR8}
Suppose the original \texttt{BPrime\_s} holds for
$a(s)=R(s)^{\rm mat}$ and the specified $K$ with
\texttt{IsQi K}.  Then ADK6 applies to
$\operatorname{Ad}_R$ on the full $M_3$, with its original
coefficient and local-condition scope.
\end{proposition}
\begin{proof}
By definition \texttt{IsQi K} says $[K:\mathbb Q]=2$ and
$K$ contains a square root of $-1$.  The algebra structure
$K\hookrightarrow\Omega$ provided by
\texttt{BPrime\_s.aut} therefore has image $\mathbb Q(i)$.
The image of restriction $G_K\to G_{\mathbb Q}$ is exactly
$\ker\eta_4$.  For example, the action on $i$ is its action on
$\mu_4$, and the nontrivial element of
$\operatorname{Gal}(\mathbb Q(i)/\mathbb Q)$ has character
value $-1$.

Unfolding \texttt{Af.AutOnAlt K iota (-1) a} gives a
representation $r_0:G_K\to\operatorname{GL}_3
(\overline{\mathbb Q}_\ell)$ whose matrix is literally
$a(s)$ on $G_K$, and says that
\[
                         r_K=\epsilon^{-1}r_0
 \tag{HSR.23}
\]
satisfies \texttt{IsAutCsdA}.  Its first conjunct is
\texttt{AutCsd0}.  Inspecting that definition gives the very same
existential cuspidal representation, regular algebraicity,
conjugate self-duality, and \texttt{AttC} condition as
\texttt{IsAutCsd}.  The normalized Satake and infinite-place
dictionary, including the geometric Frobenius convention, is
already checked in ADK3.  Thus (HSR.23) satisfies the literal
automorphy hypothesis of ADK6.  HSR1 gives continuity of $r_0$
and hence of $r_K$; the entries of the inverse are continuous as
well, so this is continuity into $\operatorname{GL}_3$.

The field $\Omega$ is an algebraic closure of this specified
copy of $K$: it is algebraically closed and algebraic over
$\mathbb Q\subset K$.  Under the above restriction map,
$\ker\epsilon|_{G_K}$ identifies with
$\ker\eta_4\cap\ker\epsilon$ in $G_{\mathbb Q}$.
The source uses $\texttt{cycZ}$ for the latter character and
$\texttt{cycBar}$ for its inclusion in
$\overline{\mathbb Q}_\ell^\times$; the inclusion is injective,
so their kernels coincide.  The \texttt{BPrime\_s.cyc}
assertion says that $r_0$ is irreducible on this exact subgroup.
The twist in (HSR.23) is one there.  Therefore $r_K$ is
irreducible on its full infinite cyclotomic kernel, exactly the
hypothesis of ADK6.  No residual or finite-level version is being
substituted.  The additional \texttt{BPrime\_s.opens} hypothesis
is retained but is unnecessary for this implication.

It remains to check the pairing and the disconnected action.
HSR1 and \texttt{IsHostFrame.cc} give
$R(c_\infty)=(1,-1)$, with $c_\infty$ induced by an actual
complex embedding.  The semidirect-product law implies, for
$s\in G_K$,
\[
 r_0(c_\infty s c_\infty^{-1})
      =Jr_0(s)^{-t}J,\qquad
 r_0(s)^tJr_0(c_\infty s c_\infty^{-1})=J.
\]
The cyclotomic character is invariant under conjugation, so
(HSR.23) gives
\[
 r_K(s)^tJr_K(c_\infty s c_\infty^{-1})
                       =\epsilon(s)^{-2}J
                       =\epsilon(s)^{1-3}J.
 \tag{HSR.24}
\]
Also $J^t=J=J^{-1}$ and, for every $X\in M_3$,
\[
 \operatorname{Ad}_R(s)X=r_K(s)Xr_K(s)^{-1},\qquad
 \operatorname{Ad}_R(c_\infty)X=-J^{-1}X^tJ.
 \tag{HSR.25}
\]
Equations (HSR.24)--(HSR.25), together with the actual complex
embedding, are every clause of \texttt{IsAdOf}, with $P=J$.
In particular the scalar line is present and complex conjugation
acts on it by $-1$.  There is no passage to trace zero and no
undetermined scalar in the pairing.  ADK6 now proves
\[
 H^1_f(\mathbb Q,\operatorname{Ad}_R)=0
\]
in the literal source sense \texttt{BlochKatoSelmerVanishes}.
\end{proof}
% END HSR8

% BEGIN HSR9
\begin{theorem}[HSR9: the original HostRigid\_s statement]
\label{claim:HSR9}
For every choice of data quantified in the source's
\texttt{HostRigid\_s}, an \texttt{IsCharZeroPt} satisfying
\texttt{BPrime\_s} is rigid in the exact sense (HSR.11).
Consequently \texttt{HOST\_RIGID\_S} is supplied by the
human proofs ADK1--ADK6, HSR1--HSR8, and the previously installed
local inputs BX1--BX4, DX1, NR19, and RX2.
\end{theorem}
\begin{proof}
Take an arbitrary continuous cocycle and arbitrary allowed
deformation from the definition of rigidity.  HSR7 makes
$Z=z\pi$ a literal Bloch--Kato cocycle for
$\operatorname{Ad}_{r\pi}$.  HSR8 and ADK6 give a single
$v\in M_3(\overline{\mathbb Q}_\ell)$ such that
\[
 z(\pi(s))=\operatorname{Ad}_r(\pi(s))v-v
                       \qquad(s\in G_{\mathbb Q}).
\]
The map $\pi$ is onto.  Set $X=-v$ and choose a preimage of
each $g\in\Gamma$ to obtain (HSR.11).  This proves every
quantifier of the original rigidity conclusion.  In particular
the primitive is one global matrix, not a family of unrelated
local primitives.
\end{proof}
% END HSR9

\subsection{Scope and source comparisons}

The proof has supplied the entire original
\texttt{HostRigid\_s} implication.  It retains the original
marked-element, parity, coefficient, valuation, and
\texttt{BPrime\_s} hypotheses, and it works at every odd
$\ell$ permitted by \texttt{HostData}.  It uses only the
$p=\ell$ instance of the numerical inequality in (HSR.2), and
only the \texttt{cyc} and \texttt{aut} parts of
\texttt{BPrime\_s}.  These observations simplify this proof;
they do not manufacture those hypotheses for a host or remove
them from any downstream source statement.

The local calculation is an explicit application of the already
proved RX2, with its finite local duality and exact crystalline
period model.  The global calculation is an explicit application
of the published theorem matched in ADK1--ADK6.  There is no new
automorphy theorem or new vanishing conjecture hidden in this
step.  Conversely this section does not supply
\texttt{HostNumber\_s}, the existence or automorphy of its points,
the quantitative congruence estimates, or the other remaining
arithmetic producers.  It closes a proof-reconstruction leaf;
it does not by itself complete the reconstruction of RH or GRH.

\begin{center}
\begin{tabular}{p{0.31\textwidth}p{0.61\textwidth}}
Original interface & Human proof and exact comparison \\
\hline
\texttt{IsCharZeroPt}, \texttt{RigDesc}
 & HSR1: norm continuity from the actual finite DVR and Baire descent.\\
\texttt{DyadicCharacterPowerNeOne}, \texttt{rig\_sub}
 & HSR2--HSR3: dyadic conductor, finite characters, stored local group,
   and the stronger witness $\epsilon(g)=1$.\\
\texttt{RigSep}, \texttt{rig\_eig}
 & HSR4: reduction of the actual characteristic polynomial in every
   triangular basis.\\
\texttt{rig\_fin}, \texttt{rig\_loc}
 & HSR5 and RX2: finite coefficients, dual-number ordinary
   crystalline calculation, and descent with the diagonal period action.\\
\texttt{rig\_diag}, \texttt{rig\_away}, \texttt{rig\_above}
 & HSR6--HSR7: explicit inertia primitives and the literal
   \texttt{decJ}/\texttt{finiteCocycles} match.\\
\texttt{BPrime\_s}, \texttt{AutOnAlt}, \texttt{IsAdOf}
 & HSR8: exact twist $\epsilon^{-1}$, full infinite cyclotomic
   kernel, and actual matrix $P=J$.\\
\texttt{HostRigid\_s}
 & HSR9: ADK6 applied to the literal global cocycle, quotient
   surjectivity, and the required sign of \texttt{cob}.\\
\end{tabular}
\end{center}

The accompanying \texttt{provenance.json} gives pinned source
paths, declaration ranges, range hashes, inherited human-proof
locations, and claim dependencies for this table.
